\chapter{Formal Status}

The public theorem in \texttt{Ringel/Statement.lean} has the form
\[
\texttt{CaseABSourceStatement} \to \texttt{ringel\_conjecture\_large}.
\]
The source package records the Case A and Case B statements needed by the proof.
The deterministic assembly and the Case C branch are formalized in Lean. The unconditional
probability, repletion, exact-colour matching, and path-absorption arguments for Cases A and B
remain the principal source-level boundary.

\chapter{Statement}

\begin{definition}\label{def:ringel_conjecture}
\lean{ringel_conjecture}
Ringel's conjecture for all $n$. This remains open.
\end{definition}

\begin{definition}\label{def:ringel_conjecture_large}
\lean{ringel_conjecture_large}
The large-$n$ theorem, conditional on \lean{CaseABSourceStatement}.
\end{definition}

\chapter{Combinatorial Primitives (Primitives.lean)}

\begin{definition}\label{def:ndColouring}
\lean{ndColouring}\leanok
Near-difference rainbow colouring of $K_{2n+1}$.
\end{definition}

\begin{theorem}\label{thm:ndColouring_isTwoFactorization}
\lean{ndColouring_isTwoFactorization}\leanok
ND-colouring is a 2-factorization.
\end{theorem}

\begin{lemma}\label{lem:ndColouring_step}
\lean{ndColouring_step}\leanok
A step by $\pm(c+1)$ realizes colour $c$.
\end{lemma}

\begin{definition}\label{def:case_predicates}
\lean{IsCaseA, IsCaseB, IsCaseC}\leanok
The three structural case predicates.
\end{definition}

\chapter{Tree Structure and Case Division}

\begin{theorem}\label{thm:split}
\lean{split}\leanok
A tree splits into many leaves or many vertex-disjoint bare paths.
\end{theorem}

\begin{theorem}\label{thm:tree_split_via_split}
\lean{tree_split_via_split}\leanok
Connects the split construction to the Case A/B/C division.
\end{theorem}

\begin{theorem}\label{thm:case_division}
\lean{case_division}\leanok
Routes an admissible tree into Case A, Case B, or Case C.
\uses{thm:split, thm:tree_split_via_split, def:case_predicates}
\end{theorem}

\chapter{Case A Source Interface (CaseA.lean and CaseASource.lean)}

\begin{definition}\label{def:walkSum}
\lean{walkSum}\leanok
Signed walk sum in the cyclic vertex set.
\end{definition}

\begin{lemma}\label{lem:tree_embed}
\lean{tree_embed}\leanok
Embeds a tree through signed edge increments.
\end{lemma}

\begin{definition}\label{def:valid_caseA_embedding}
\lean{valid_caseA_embedding}\leanok
Records the core embedding, leaf positions, injectivity, and rainbow conditions needed by the
Case A finishing step.
\end{definition}

\begin{structure}\label{def:caseA_joint_output}
\lean{CaseAJointOutput}\leanok
Records a finite Case A output together with its core map, reservoirs, and perfect rainbow
matching.
\end{structure}

\begin{theorem}\label{thm:caseA_source}
\lean{CaseASourceStatement}
The eventual Case A source statement used by the public source package.
\end{theorem}

\chapter{Case B Source Interface (CaseB.lean and CaseSource.lean)}

\begin{definition}\label{def:caseB_source}
\lean{CaseBSourceStatement}
The eventual Case B source statement used by the public source package.
\end{definition}

The fixed-endpoint path construction, its random availability estimates, and exact collective
colour use remain to be formalized from the MPS source.

\chapter{Case C --- Small Core (CaseC.lean)}

\begin{lemma}\label{lem:caseC_embedding_exists}
\lean{caseC_embedding_exists}\leanok
Constructs the Case C rainbow embedding for sufficiently large $n$.
\end{lemma}

\begin{theorem}\label{thm:caseC_rainbow}
\lean{caseC_rainbow}\leanok
Produces a rainbow copy in Case C.
\uses{lem:caseC_embedding_exists}
\end{theorem}

\chapter{Probability and Reservoir Infrastructure}

\begin{definition}\label{def:prob_event}
\lean{prob_event}\leanok
Finite probability space used by the formal development.
\end{definition}

\begin{lemma}\label{lem:exists_of_prob_gt_zero}
\lean{exists_of_prob_gt_zero}\leanok
Positive probability implies existence.
\end{lemma}

The near-embedding and reservoir modules formalize finite sample spaces, concentration
bookkeeping, and retained sets. They do not by themselves prove the full MPS Case A/B source
statements.

\chapter{Proof Assembly (Spine.lean)}

\begin{theorem}\label{thm:decomp_of_rainbow_copy}
\lean{decomp_of_rainbow_copy}\leanok
The Kotzig cyclic-shift construction turns one rainbow copy into an edge decomposition.
\end{theorem}

\begin{theorem}\label{thm:rainbow_copy_exists}
\lean{rainbow_copy_exists}
\uses{thm:case_division, thm:caseA_source, def:caseB_source, thm:caseC_rainbow}
The source package yields an eventual rainbow copy.
\end{theorem}

\begin{theorem}\label{thm:ringel_conjecture_large_via_spine}
\lean{ringel_conjecture_large_via_spine}
\uses{thm:decomp_of_rainbow_copy, thm:rainbow_copy_exists}
Assembles the large-$n$ theorem from the source package.
\end{theorem}
